\clearpage
\VerdierPhysicalPageAnchor{0025}
这里 $\mathrm{RH}(X,{\ })$ 表示 $X$ 上“整体截面”函子的全右导出函子
（$\mathrm{RH}(Y,{\ })$ 与 $\mathrm{RH}(X\times Y,{\ })$ 同理），而
$\overset{\mathrm L}{\otimes}_A$ 表示张量积函子的全左导出函子。
（假定 $X$ 与 $Y$ 的上同调维数有限，例如它们是有限胞腔复形。）
$\text{(Kü)}_{\mathrm{Der}}$ 两边都是在拟同构意义下确定的 $A$-模复形，
其上同调模分别是 $\text{(Kü)}_{\mathrm{coh}}^{4}$ 中第一与第二个谱序列
的收敛目标。当然，当 $F$ 与 $G$ 是 $X$ 与 $Y$ 上适当有界的层复形时，
公式 $\text{(Kü)}_{\mathrm{Der}}$ 仍然成立。它在实际运算中完全可用，并且
在有多个因子时，也能表述交换性与结合性。

\VerdierSemanticAnchor{intro-5-base-change}
标量扩张在导出范畴中给出类似公式。若 $A\to B$ 是一个 $A$-代数，$F$ 是
$X$ 上的 $A$-模层，则有关系
\[
\mathrm{RH}(X,F)\overset{\mathrm L}{\otimes}_A B
\xrightarrow{\sim}
\mathrm{RH}\bigl(X,F\overset{\mathrm L}{\otimes}_A B\bigr).
\]

\VerdierSemanticAnchor{intro-5-naive-formulas}
形象地说，一般而言，\emph{传统上不成立的朴素公式，一旦在导出范畴中工作
便会成立}。

\VerdierSemanticAnchor{intro-5-l-adic}
系数环 $\mathbb Z/\ell^r\mathbb Z$ 的重要性尤其来自概形的 étale 上同调。
已知 [6] 在该理论中必须主要处理挠层。取
$\mathbb Z/\ell^r\mathbb Z$ 系数上同调的逆极限，就得到 $\ell$-进上同调，
即取值于 $\ell$-进整数环 $\widehat{\mathbb Z}_{\ell}$ 的上同调。
公式 $\text{(Kü)}_{\mathrm{Der}}$ 在这一背景下仍成立（例如 $X$ 与 $Y$
是代数闭域上的完备概形，且 $A=\mathbb Z/\ell^r\mathbb Z$），并可立即
通过逆极限，给出 $\widehat{\mathbb Z}_{\ell}$-层的类似公式。又因为
$\widehat{\mathbb Z}_{\ell}$ 是主理想整环，公式
$\text{(Kü)}_{\mathrm{Der}}$ 在取上同调后给出
$\text{(Kü)}_{\mathrm{coh}}^{3}$ 型公式；再与 $\mathbb Q_{\ell}$ 作
张量积，则得到 $\mathbb Q_{\ell}$-层的
$\text{(Kü)}_{\mathrm{coh}}^{1}$ 型公式。

\VerdierSemanticAnchor{intro-6}
\noindent\textbf{6.}\quad
本文系统展开上面概述的这些思想。全书安排如下。

\VerdierSemanticAnchor{intro-6-plan-chapter-one}
第一章研究加性范畴的复形范畴。这些范畴带有分次。更确切地说，它们配备
平移函子（沿给定方向平移次数，并相应改变微分的符号）。
